\subsection{Torsion and Quotients}

Throughout this subsection, $R$ denotes a nonzero commutative ring with identity and $M$ denotes an $R$-module.

\begin{definition} \label{def:module-torsion} \hlt{Torsion}
\begin{enumerate}[label=(\roman*)]
\setlength{\itemsep}{0pt}
\item A nonzero element $v\in M$ is a \hlt{torsion element} if $\exists\,0\neq r\in R$ such that $rv=0$.
\item $M$ is \hlt{torsion-free} if it has no torsion elements.
\item $M$ is a \hlt{torsion module} if every nonzero element of $M$ is torsion.
\item The \hlt{torsion set} of $M$ is $M_{\operatorname{tor}}\defeq\{0\}\cup\{v\in M:v$ is torsion$\}$
\end{enumerate}
\end{definition}

\begin{example} \label{ex:zn-is-torsion-module} \hlt{$\mathbb{Z}/n\mathbb{Z}$ as a Torsion Module}\\
For $n>1$, $\mathbb{Z}/n\mathbb{Z}$ is a $\mathbb{Z}$-module via $k\cdot \overline{a}=\overline{ka}$. For every $\overline{a}\in\mathbb{Z}/n\mathbb{Z}$, $n\overline{a}=\overline{0}$.\\
Hence $\mathbb{Z}/n\mathbb{Z}$ is a torsion module, and $\{\overline{a}\}$ is linearly dependent whenever $\overline{a}\neq\overline{0}$.
\end{example}

\begin{proposition} \label{prop:torsion-submodule-over-domain} \hlt{Torsion Submodule over an Integral Domain}\\
If $R$ is an integral domain, then $M_{\operatorname{tor}}\leq M$, $(M/M_{\operatorname{tor}})_{\operatorname{tor}}=\{0\}$.
\begin{proof}
Let $x,y\in M_{\operatorname{tor}}$ and $r,s\in R$.\\
Choose $a,b\in R\setminus\{0\}$ with $ax=0$ and $by=0$, taking $a=1$ if $x=0$ and $b=1$ if $y=0$.\\
Since $R$ is a domain, $ab\neq0$, and $ab(rx+sy)=rb(ax)+sa(by)=0$. Thus $rx+sy\in M_{\operatorname{tor}}$, so $M_{\operatorname{tor}}\leq M$.\\
Now let $0\neq r\in R$ and suppose $r(m+M_{\operatorname{tor}})=M_{\operatorname{tor}}$. Then $rm\in M_{\operatorname{tor}}$, so $srm=0$ for some $0\neq s\in R$; if $rm=0$, take $s=1$. Since $sr\neq0$, $m\in M_{\operatorname{tor}}$, hence $m+M_{\operatorname{tor}}=M_{\operatorname{tor}}$.\\
\hlt{Structural Commentary}: Domains let annihilators be multiplied without becoming zero.\\
\hlt{Key Insight}: Domains make torsion closed and quotients torsion-free.
\end{proof}
\end{proposition}

\begin{definition} \label{def:module-annihilator} \hlt{Annihilators}
\begin{enumerate}[label=(\roman*)]
\setlength{\itemsep}{0pt}
\item For $v\in M$, the \hlt{annihilator} of $v$ is $\operatorname{Ann}_R(v)\defeq\{r\in R:rv=0\}$.
\item For $N\leq M$, the \hlt{annihilator} of $N$ is $\operatorname{Ann}_R(N)\defeq\{r\in R:rN=\{0\}\}$.
\end{enumerate}
\end{definition}

\begin{proposition} \label{prop:annihilator-basic-properties} \hlt{Annihilator Properties}\\
For $v\in M$ and $N\leq M$:
\begin{enumerate}[label=(\roman*)]
\setlength{\itemsep}{0pt}
\item $\operatorname{Ann}_R(v)\trianglelefteq R$ and $\operatorname{Ann}_R(N)\trianglelefteq R$.
\item $v$ is torsion $\Longleftrightarrow$ $v\neq0$ and $\operatorname{Ann}_R(v)\neq\{0\}$.
\item $A\leq B\leq M \Longrightarrow \operatorname{Ann}_R(B)\subseteq\operatorname{Ann}_R(A)$.
\end{enumerate}
\begin{proof}
\begin{enumerate}[label=\roman*.]
\setlength{\itemsep}{0pt}
\item[$(i)$] Since $0\in\operatorname{Ann}_R(v)$ and, for $a,b\in\operatorname{Ann}_R(v)$ and $r,s\in R$, $(ra+sb)v=r(av)+s(bv)=0$, $\operatorname{Ann}_R(v)\trianglelefteq R$. The same computation for every $n\in N$ gives $\operatorname{Ann}_R(N)\trianglelefteq R$.
\item[$(ii)$]$(\Rightarrow)$ If $v$ is torsion, then $v\neq0$ and $\exists\,0\neq r\in R$ with $rv=0$, so $\operatorname{Ann}_R(v)\neq\{0\}$.
\item[$(ii)$]$(\Leftarrow)$ If $v\neq0$ and $\operatorname{Ann}_R(v)\neq\{0\}$, choose $0\neq r\in\operatorname{Ann}_R(v)$. Then $rv=0$, so $v$ is torsion.
\item[$(iii)$] If $A\leq B\leq M$ and $r\in\operatorname{Ann}_R(B)$, then $rA=\{0\}$; hence $r\in\operatorname{Ann}_R(A)$.
\end{enumerate}
\hlt{Structural Commentary}: Annihilators reverse inclusion.\\
\hlt{Key Insight}: Killing more vectors gives fewer scalars.
\end{proof}
\end{proposition}

\begin{proposition} \label{prop:finitely-generated-torsion-module-annihilator} \hlt{Torsion Generators Give a Nonzero Annihilator over a Domain}\\
Assume $R$ is an integral domain. If $M=\langle\!\langle u_1,\ldots,u_n\rangle\!\rangle_R$, $u_i\in M_{\operatorname{tor}}\ \forall i$, then $\operatorname{Ann}_R(M)\neq\{0\}$.
\begin{proof}
Choose $0\neq a_i\in R$ with $a_iu_i=0$, and set $a\defeq\prod\limits_{i=1}^{n}a_i$. Since $R$ is a domain, $a\neq0$.\\
For each $i$, $au_i=0$, hence $aM=\{0\}$ because the $u_i$ generate $M$. Thus $0\neq a\in\operatorname{Ann}_R(M)$.\\
\hlt{Structural Commentary}: Finite generation lets many annihilators be compressed into one product.\\
\hlt{Key Insight}: A finite torsion generating set gives a global order for the module.
\end{proof}
\end{proposition}

\begin{definition} \label{def:module-homomorphism} \hlt{Module Homomorphism}\\
Let $M,N$ be $R$-modules and let $\tau:M\to N$.
\begin{enumerate}[label=(\roman*)]
\setlength{\itemsep}{0pt}
\item $\tau$ is \hlt{$R$-linear} if, $\forall r,s\in R$ and $u,v\in M$, $\tau(ru+sv)=r\tau(u)+s\tau(v)$, and is then called an \hlt{$R$-homomorphism}.
\item The set of all $R$-homomorphisms $M\to N$ is $\operatorname{Hom}_R(M,N)$.
\item An $R$-homomorphism $M\to M$ is an \hlt{endomorphism}.
\item An injective, surjective, or bijective $R$-homomorphism is a \hlt{monomorphism}, \hlt{epimorphism}, or \hlt{isomorphism}.
\end{enumerate}
\end{definition}

\begin{proposition} \label{prop:hom-module-kernel-image} \hlt{Hom, Kernel, and Image}\\
Let $M,N$ be $R$-modules and $\tau\in\operatorname{Hom}_R(M,N)$.
\begin{enumerate}[label=(\roman*)]
\setlength{\itemsep}{0pt}
\item With $(a\sigma+b\rho)(m)\defeq a\sigma(m)+b\rho(m)$, $\operatorname{Hom}_R(M,N)$ is an $R$-module.
\item $\ker\tau\defeq\tau^{-1}(\{0\})\leq M$.
\item $\operatorname{im}\tau\defeq\tau(M)\leq N$.
\item $\tau$ is a monomorphism $\Longleftrightarrow\ker\tau=\{0\}$.
\end{enumerate}
\begin{proof} By Theorem \ref{thm:submodule-test}, for $A\subseteq P$ in an $R$-module $P$, $A\leq P$ $\Longleftrightarrow$ $A\neq\emptyset$ and $\forall a,b\in R$, $\forall x,y\in A$, $ax+by\in A$.
\begin{enumerate}[label=(\roman*)]
\setlength{\itemsep}{0pt}
\item For $\sigma,\rho\in\operatorname{Hom}_R(M,N)$ and $a,b,r,s\in R$, $(a\sigma+b\rho)(ru+sv) =r(a\sigma+b\rho)(u)+s(a\sigma+b\rho)(v)$, where $u,v\in M$, using commutativity of $R$. Since $0\in\operatorname{Hom}_R(M,N)$, the displayed criterion gives the claim inside the pointwise module of maps $M\to N$.
\item Since $0\in\ker\tau$, and $x,y\in\ker\tau$, $a,b\in R$ imply $\tau(ax+by)=a\tau(x)+b\tau(y)=0$,\\
so the displayed criterion gives $\ker\tau\leq M$.
\item Since $0=\tau(0)\in\operatorname{im}\tau$, and $y_i=\tau(x_i)\in\operatorname{im}\tau$ $(i=1,2)$, $a,b\in R$ imply $ay_1+by_2=\tau(ax_1+bx_2)\in\operatorname{im}\tau$, so the displayed criterion gives $\operatorname{im}\tau\leq N$.
\item If $\tau$ is injective, then $\tau(m)=0=\tau(0)$ implies $m=0$, so $\ker\tau=\{0\}$.\\
Conversely, if $\ker\tau=\{0\}$ and $\tau(m)=\tau(m')$, then $\tau(m-m')=0$, hence $m=m'$.
\end{enumerate}
\hlt{Structural Commentary}: Kernels and images are the canonical submodules attached to a module map.\\
\hlt{Key Insight}: Homomorphisms preserve precisely the finite $R$-linear combinations that define submodules.
\end{proof}
\end{proposition}

\begin{definition} \label{def:quotient-module} \hlt{Quotient Module}\\
Let $S\leq M$. Define $u\equiv v\pmod S \Longleftrightarrow u-v\in S$.\\
The \hlt{quotient module} is $M/S\defeq\{u+S:u\in M\}$; $u+S\defeq\{u+s:s\in S\}$,\\
with operations $(u+S)+(v+S)\defeq(u+v)+S$; $r(u+S)\defeq ru+S$.
\end{definition}

\begin{proposition} \label{prop:quotient-module-well-defined} \hlt{Quotient Operations Are Well-Defined}\\
Let $S\leq M$. The operations of Definition \ref{def:quotient-module} are well-defined and make $M/S$ an $R$-module.\\
Moreover, $0_{M/S}=S$, $\pi:M\twoheadrightarrow M/S$, $\pi(u)\defeq u+S$, $\ker\pi=S.
$.
\begin{proof}
If $u\equiv u'\pmod S$ and $v\equiv v'\pmod S$, then $(u+v)-(u'+v')=(u-u')+(v-v')\in S$\\
and $ru-ru'=r(u-u')\in S$. Thus $+$ and $R$-scaling on cosets are representative-independent. The module axioms descend from $M$. Finally, $0_{M/S}=0+S=S$, $\pi$ is surjective, and $\ker\pi=\{u\in M:u+S=S\}=\{u\in M:u\in S\}=S$.
\hlt{Structural Commentary}: Quotients collapse exactly the submodule $S$ to zero.\\
\hlt{Key Insight}: Two representatives define the same coset precisely when their difference is invisible modulo $S$.
\end{proof}
\end{proposition}

\begin{example} \label{ex:quotient-free-module-need-not-free} \hlt{A Quotient of a Free Module Need Not Be Free}\\
As a $\mathbb{Z}$-module, $\mathbb{Z}$ is free on $\{1\}$. For $n>1$, the submodule $n\mathbb{Z}$ is cyclic, but the quotient $\mathbb{Z}/n\mathbb{Z}$ is torsion by Example \ref{ex:zn-is-torsion-module}. Hence $\mathbb{Z}/n\mathbb{Z}$ is not free as a $\mathbb{Z}$-module.
\end{example}

\begin{theorem} \label{thm:module-correspondence-theorem} \hlt{Correspondence Theorem for Modules}\\
Let $S\leq M$ and let $\pi:M\twoheadrightarrow M/S$ be the quotient map. Then $\Phi:\{T\leq M:S\leq T\}\longrightarrow\{U\leq M/S\}$, where\\
$\Phi(T)\defeq T/S$ is an order isomorphism, with inverse $\Psi(U)\defeq\pi^{-1}(U)$.
\begin{proof} If $S\leq T\leq M$, then $\Phi(T)=\pi(T)\leq M/S$.\\
If $U\leq M/S$, then $\Psi(U)=\pi^{-1}(U)\leq M$ and $S=\ker\pi\leq\Psi(U)$.\\
Moreover, $\Psi\Phi(T)=\pi^{-1}(T/S)=T$ and $\Phi\Psi(U)=\pi(\pi^{-1}(U))=U$, where the second equality uses surjectivity of $\pi$. Finally, $T_1\leq T_2\Longrightarrow T_1/S\leq T_2/S$, and $U_1\leq U_2\Longrightarrow \pi^{-1}(U_1)\leq \pi^{-1}(U_2)$. Thus $\Phi$ is an order isomorphism.\\
\hlt{Structural Commentary}: Submodules above $S$ are exactly the submodules visible after $S$ is collapsed.\\
\hlt{Key Insight}: Pull back along the quotient map to recover the original submodule.
\end{proof}
\end{theorem}

\begin{theorem} \label{thm:module-isomorphism-theorems} \hlt{Module Isomorphism Theorems}\\
The following canonical maps are $R$-module isomorphisms.
\begin{enumerate}[label=(\roman*)]
\setlength{\itemsep}{0pt}
\item If $\tau\in\operatorname{Hom}_R(M,N)$, then $\overline{\tau}:M/\ker\tau\xrightarrow{\sim}\operatorname{im}\tau$ and $\overline{\tau}(m+\ker\tau)\defeq\tau(m)$.
\item If $S,T\leq M$, then $\theta:S/(S\cap T)\xrightarrow{\sim}(S+T)/T$ and $\theta\bigl(s+(S\cap T)\bigr)\defeq s+T$.
\item If $S\leq T\leq M$, then $\eta:(M/S)/(T/S)\xrightarrow{\sim}M/T$ and $\eta\bigl((m+S)+T/S\bigr)\defeq m+T$.
\end{enumerate}
\begin{proof}
\begin{enumerate}[label=(\roman*)]
\setlength{\itemsep}{0pt}
\item If $m-m'\in\ker\tau$, then $\tau(m)=\tau(m')$, so $\overline{\tau}$ is well-defined and $R$-linear.\\
Also $\operatorname{im}\overline{\tau}=\operatorname{im}\tau$ and $\ker\overline{\tau}=\{m+\ker\tau:\tau(m)=0\}=\{0_{M/\ker\tau}\}$.
\item Apply (i) to $\varphi\in\operatorname{Hom}_R(S,(S+T)/T)$, $\varphi(s)\defeq s+T$.\\
Then $\operatorname{im}\varphi=(S+T)/T$ and $\ker\varphi=\{s\in S:s\in T\}=S\cap T$.
\item Apply (i) to $\psi\in\operatorname{Hom}_R(M/S,M/T)$, $\psi(m+S)\defeq m+T$.\\
Then $\operatorname{im}\psi=M/T$ and $\ker\psi=\{m+S:m\in T\}=T/S$.
\end{enumerate}
\hlt{Structural Commentary}: Each theorem identifies a quotient by the exact ambiguity in a natural map.\\
\hlt{Key Insight}: Image is domain modulo kernel, repeatedly.
\end{proof}
\end{theorem}

\begin{definition} \label{def:module-direct-sums} \hlt{Direct Sums and Direct Summands}
\begin{enumerate}[label=(\roman*)]
\setlength{\itemsep}{0pt}
\item For $R$-modules $M_1,\ldots,M_n$, their \hlt{external direct sum} is $\bigoplus\limits_{i=1}^{n}M_i\defeq\{(m_i)_{i=1}^{n}:m_i\in M_i\}$\\
with componentwise $R$-module operations.
\item For submodules $\{S_i:i\in I\}$ of $M$, their \hlt{internal direct sum} is denoted $M=\bigoplus_{i\in I}S_i$,\\
which means $M=\sum\limits_{i\in I}S_i$ and $\forall i\in I$, $S_i\cap\sum_{j\neq i}S_j=\{0\}$.
\item A submodule $S\leq M$ is a \hlt{direct summand} if $\exists\,T\leq M$ such that $M=S\oplus T$.\\
Such a $T$ is a \hlt{complement} of $S$ in $M$.
\end{enumerate}
\end{definition}

\begin{proposition} \label{prop:module-direct-sum-characterizations} \hlt{Direct Sum Characterizations}\\
For a family $\{S_i:i\in I\}$ of distinct submodules of $M$, the following conditions are equivalent:
\begin{enumerate}[label=(\roman*)]
\setlength{\itemsep}{0pt}
\item $\sum\limits_{i\in I}S_i$ is direct, i.e. $\forall k\in I,\ S_k\cap\sum\limits_{i\neq k}S_i=\{0\}$.
\item For every $F\subseteq_{\operatorname{fin}} I$ and $x_i\in S_i$ $(i\in F)$, $\sum\limits_{i\in F}x_i=0$ $\Longrightarrow$ $x_i=0$ $\forall i \in F$.
\item Every $0\neq x\in\sum_{i\in I}S_i$ has a unique finite nonzero decomposition,\\
$\exists!\,(F,(x_i)_{i\in F})$ where $F\subseteq_{\operatorname{fin}} I$; $0\neq x_i\in S_i$; $x=\sum\limits_{i\in F}x_i$.
\end{enumerate}
\begin{proof}
\begin{enumerate}[label=(\roman*).]
\setlength{\itemsep}{0pt}
\item[$(i)$]$\Rightarrow (ii)$. If (i) holds and $\sum\limits_{i\in F}x_i=0$, then any $x_k\neq0$ gives $x_k=-\sum\limits_{i\in F\setminus\{k\}}x_i\in S_k\cap\sum\limits_{i\neq k}S_i$, a contradiction.
\item[$(ii)$]$\Rightarrow (i)$. If $x\in S_k\cap\sum\limits_{i\neq k}S_i$, write $x=\sum\limits_{i\in F}x_i$ with $F\subseteq_{\operatorname{fin}} I\setminus\{k\}$ and $x_i\in S_i$.\\
Then $x-\sum\limits_{i\in F}x_i=0$, so (ii) gives $x=0$.
\item[$(ii)$]$\Rightarrow (iii)$. Existence follows from the definition of $\sum_i S_i$ after deleting zero terms.\\
If $\sum\limits_{i\in F}x_i=\sum\limits_{j\in G}y_j$, $0\neq x_i\in S_i$, $0\neq y_j\in S_j$, then, over $H=F\cup G$, missing terms set to $0$, $\sum\limits_{h\in H}(x_h-y_h)=0$.\\
By (ii), $x_h=y_h$ for all $h\in H$, hence $F=G$ and the terms agree.
\item[$(iii)$]$\Rightarrow (ii)$. If $\sum_{i\in F}x_i=0$, set $F_0\defeq\{i\in F:x_i\neq0\}$.\\
If $F_0\neq\emptyset$, choose $k\in F_0$. If $F_0=\{k\}$, then $x_k=0$, contradiction; otherwise $0\neq x_k=-\sum\limits_{i\in F_0\setminus\{k\}}x_i$,\\
giving two finite nonzero decompositions of $x_k$, contradiction. Hence $F_0=\emptyset$.
\end{enumerate}
\hlt{Structural Commentary}: Directness is the absence of hidden cancellation among different summands.\\
\hlt{Key Insight}: Unique decomposition is controlled by the zero vector.
\end{proof}
\end{proposition}

\begin{example} \label{ex:z-submodules-mostly-not-complemented} \hlt{Most Submodules of $\mathbb{Z}$ Are Not Complemented}\\
As a $\mathbb{Z}$-module, every submodule of $\mathbb{Z}$ is $n\mathbb{Z}$ for a unique $n\geq0$. For $m,n>0$, $m\mathbb{Z}\cap n\mathbb{Z}=\operatorname{lcm}(m,n)\mathbb{Z}\neq\{0\}$.\\
Thus $m\mathbb{Z}\oplus n\mathbb{Z}\neq\mathbb{Z}$ for all $m,n>0$, so the only complemented submodules of $\mathbb{Z}$ are $\{0\}$ and $\mathbb{Z}$.
\end{example}

\begin{theorem} \label{thm:complements-and-quotients-modules} \hlt{Complements and Quotients}\\
If $M=S\oplus T$, then $\pi|_T:T\xrightarrow{\sim}M/S$, $t\mapsto t+S$, where $\pi:M\twoheadrightarrow M/S$ is the quotient map.\\
Hence any two complements of $S$ are isomorphic.
\begin{proof} $\ker(\pi|_T)=T\cap S=\{0\}$. Also, for $m=s+t$ with $s\in S$, $t\in T$, one has $m+S=t+S$; hence $\operatorname{im}(\pi|_T)=M/S$. Thus $\pi|_T$ is an isomorphism, and all complements are isomorphic to $M/S$.\\
\hlt{Structural Commentary}: A complement gives a concrete copy of the quotient inside the original module.\\
\hlt{Key Insight}: Quotienting by $S$ remembers exactly the complementary coordinate.
\end{proof}
\end{theorem}

\begin{theorem} \label{thm:direct-summands-extend-epimorphisms} \hlt{Direct Summands Extend Epimorphisms}\\
Let $N\leq M$ and let $M_1$ be an $R$-module.
\begin{enumerate}[label=(\roman*)]
\setlength{\itemsep}{0pt}
\item If $M=N\oplus H$ and $\sigma:N\twoheadrightarrow M_1$, then $\exists$ a unique $\overline{\sigma}\in\operatorname{Hom}_R(M,M_1)$ such that\\
$\overline{\sigma}(n+h)\defeq\sigma(n)$, $n\in N$, $h\in H$, $\overline{\sigma}|_N=\sigma$, $\overline{\sigma}:M\twoheadrightarrow M_1$, and $\ker\overline{\sigma}=\ker\sigma\oplus H$.
\item If $\sigma:N\xrightarrow{\sim}M_1$, then $\mathcal{E}_{\sigma}\defeq\{\overline{\sigma}:M\twoheadrightarrow M_1:\overline{\sigma}|_N=\sigma\}$ $\longleftrightarrow$ $\mathcal{C}_N(M)\defeq\{H\leq M:M=N\oplus H\}$ via $\overline{\sigma}\mapsto\ker\overline{\sigma}$ and $H\mapsto\overline{\sigma}_H$, where $\overline{\sigma}_H(n+h)\defeq\sigma(n)$.
\end{enumerate}
\begin{proof}
\begin{enumerate}[label=(\roman*)]
\setlength{\itemsep}{0pt}
\item Direct-sum uniqueness makes $\overline{\sigma}$ well-defined, and componentwise computation gives $R$-linearity.\\
Moreover, $\operatorname{im}\overline{\sigma}=\operatorname{im}\sigma=M_1$; $\ker\overline{\sigma}=\{n+h:\sigma(n)=0\}=\ker\sigma\oplus H$. This shows uniqueness.
\item If $H\in\mathcal{C}_N(M)$, then (i) gives $\overline{\sigma}_H\in\mathcal{E}_{\sigma}$ and, since $\ker\sigma=\{0\}$, $\ker\overline{\sigma}_H=H$.\\
Conversely, if $\overline{\sigma}\in\mathcal{E}_{\sigma}$, then $N\cap\ker\overline{\sigma}=\{0\}$, and for each $m\in M$, $n\defeq\sigma^{-1}(\overline{\sigma}(m))\in N$, $m-n\in\ker\overline{\sigma}$.\\
Thus $M=N\oplus\ker\overline{\sigma}$, so the two displayed assignments are inverse.
\end{enumerate}
\hlt{Structural Commentary}: A complement supplies the directions on which an extension may vanish.\\
\hlt{Key Insight}: Extending an isomorphism is equivalent to choosing its kernel complement.
\end{proof}
\end{theorem}

\begin{definition} \label{def:module-retract} \hlt{Retract}\\
Let $N\leq M$.
\begin{enumerate}[label=(\roman*)]
\setlength{\itemsep}{0pt}
\item A \hlt{retraction} onto $N$ is a map $\rho\in\operatorname{Hom}_R(M,N)$ such that $\rho|_N=\operatorname{id}_N$.
\item The submodule $N$ is a \hlt{retract} of $M$ iff $\exists\,\rho\in\operatorname{Hom}_R(M,N)$ with $\rho|_N=\operatorname{id}_N$.
\end{enumerate}
\end{definition}

\begin{corollary} \label{cor:retract-iff-direct-summand-module} \hlt{Retracts Are Direct Summands}\\
For $N\leq M$, $N$ is a retract of $M$ $\Longleftrightarrow$ $\exists\,H\leq M$ such that $M=N\oplus H$.
\begin{proof}
\begin{enumerate}[label=\roman*.]
\setlength{\itemsep}{0pt}
\item[$(\Rightarrow)$] Let $\rho:M\to N$ satisfy $\rho|_N=\operatorname{id}_N$. For every $m\in M$, $m=\rho(m)+(m-\rho(m))$, $\rho(m)\in N$, $m-\rho(m)\in\ker\rho$.\\
Moreover $N\cap\ker\rho=\{0\}$, so $M=N\oplus\ker\rho$.
\item[$(\Leftarrow)$] If $M=N\oplus H$, define $\rho:M\to N$ by $\rho(n+h)\defeq n$, $n\in N$, $h\in H$.\\
Then $\rho\in\operatorname{Hom}_R(M,N)$ and $\rho|_N=\operatorname{id}_N$, so $N$ is a retract.
\end{enumerate}
\hlt{Structural Commentary}: Retractions are projections onto submodules.\\
\hlt{Key Insight}: A submodule is split exactly when it can be recovered by an idempotent map.
\end{proof}
\end{corollary}

\begin{definition} \label{def:one-sided-inverses-module-maps} \hlt{One-Sided Inverses}\\
Let $\tau\in\operatorname{Hom}_R(A,B)$.
\begin{enumerate}[label=(\roman*)]
\setlength{\itemsep}{0pt}
\item A \hlt{left inverse} of $\tau$ is $\lambda\in\operatorname{Hom}_R(B,A)$ such that $\lambda\tau=\operatorname{id}_A$.
\item A \hlt{right inverse} of $\tau$ is $\rho\in\operatorname{Hom}_R(B,A)$ such that $\tau\rho=\operatorname{id}_B$.
\item A \hlt{two-sided inverse} of $\tau$ is $\upsilon\in\operatorname{Hom}_R(B,A)$ such that $\upsilon\tau=\operatorname{id}_A$ and $\tau\upsilon=\operatorname{id}_B$.
\end{enumerate}
\end{definition}

\begin{theorem} \label{thm:module-one-sided-inverses-and-complements} \hlt{One-Sided Inverses and Complements}\\
Let $\sigma:M\to M_1$ be an $R$-homomorphism.\\
For $S\leq N$, set $\mathcal{C}_S(N)\defeq\{H\leq N:N=S\oplus H\}$, $\mathcal{L}_\sigma\defeq\{\lambda\in\operatorname{Hom}_R(M_1,M):\lambda\sigma=\operatorname{id}_M\}$,\\
and $\mathcal{R}_\sigma\defeq\{\rho\in\operatorname{Hom}_R(M_1,M):\sigma\rho=\operatorname{id}_{M_1}\}$.
\begin{enumerate}[label=(\roman*)]
\setlength{\itemsep}{0pt}
\item If $\sigma$ is injective, then $\mathcal{L}_\sigma\xrightarrow{\sim}\mathcal{C}_{\operatorname{im}\sigma}(M_1)$ by $\lambda\mapsto\ker\lambda$.
\item If $\sigma$ is surjective, then $\mathcal{R}_\sigma\xrightarrow{\sim}\mathcal{C}_{\ker\sigma}(M)$ by $\rho\mapsto\operatorname{im}\rho$.
\end{enumerate}
\begin{proof}
\begin{enumerate}[label=(\roman*)]
\setlength{\itemsep}{0pt}
\item For $H\in\mathcal{C}_{\operatorname{im}\sigma}(M_1)$, define $\lambda_H(\sigma m+h)\defeq m$. Directness and injectivity give well-definedness; $\lambda_H\sigma=\operatorname{id}_M$ and $\ker\lambda_H=H$. Conversely, if $\lambda\in\mathcal{L}_\sigma$, then $\operatorname{im}\sigma\cap\ker\lambda=\{0\}$, since $\sigma m\in\ker\lambda\Rightarrow m=0$, and $y=\sigma\lambda y+(y-\sigma\lambda y)$ with $y-\sigma\lambda y\in\ker\lambda$. Thus $M_1=\operatorname{im}\sigma\oplus\ker\lambda$, and the constructions are inverse.
\item For $H\in\mathcal{C}_{\ker\sigma}(M)$, $\sigma|_H:H\to M_1$ is an isomorphism; let $\rho_H=(\sigma|_H)^{-1}$ followed by $H\hookrightarrow M$.\\
Then $\rho_H\in\mathcal{R}_\sigma$ and $\operatorname{im}\rho_H=H$. Conversely, if $\rho\in\mathcal{R}_\sigma$, then $\rho$ is injective, $\operatorname{im}\rho\cap\ker\sigma=\{0\}$, and $m=(m-\rho\sigma m)+\rho\sigma m$ with $m-\rho\sigma m\in\ker\sigma$, $\rho\sigma m\in\operatorname{im}\rho$.\\
Thus $M=\ker\sigma\oplus\operatorname{im}\rho$, and the constructions are inverse.
\end{enumerate}
\hlt{Structural Commentary}: One-sided inverses exist exactly when the relevant image or kernel splits.\\
\hlt{Key Insight}: Injective maps need a complemented image; surjective maps need a complemented kernel.
\end{proof}
\end{theorem}

\begin{proposition} \label{prop:split-epimorphism-modules} \hlt{Split Epimorphisms Split the Domain}\\
Let $\sigma:M\twoheadrightarrow N$ be an $R$-homomorphism and set $\mathcal{R}_\sigma\defeq\{\rho\in\operatorname{Hom}_R(N,M):\sigma\rho=\operatorname{id}_N\}$. These are equivalent:
\begin{enumerate}[label=(\roman*)]
\setlength{\itemsep}{0pt}
\item $\mathcal{R}_\sigma\neq\emptyset$.
\item $\exists\,H\leq M$ such that $M=\ker\sigma\oplus H$.
\end{enumerate}
Moreover, $\rho\in\mathcal{R}_\sigma\Longrightarrow M=\ker\sigma\oplus\operatorname{im}\rho$.
\begin{proof}
\begin{enumerate}[label=\roman*.]
\setlength{\itemsep}{0pt}
\item[$(\Rightarrow)$] Let $\rho\in\mathcal{R}_\sigma$. If $\rho n=0$, then $n=\sigma\rho n=0$, so $\rho$ is injective.\\
If $\rho n\in\ker\sigma$, then $n=\sigma\rho n=0$, hence $\operatorname{im}\rho\cap\ker\sigma=\{0\}$.\\
For $m\in M$, $m=(m-\rho\sigma m)+\rho\sigma m$, with $m-\rho\sigma m\in\ker\sigma$ and $\rho\sigma m\in\operatorname{im}\rho$. Thus $M=\ker\sigma\oplus\operatorname{im}\rho$.
\item[$(\Leftarrow)$] Let $M=\ker\sigma\oplus H$. Then $\ker(\sigma|_H)=H\cap\ker\sigma=\{0\}$, and $\operatorname{im}(\sigma|_H)=N$ since $n=\sigma m$, $m=k+h$, gives $n=\sigma h$. Thus $\sigma|_H:H\xrightarrow{\sim}N$, and $\rho\defeq\iota_H(\sigma|_H)^{-1}\in\mathcal{R}_\sigma$, where $\iota_H:H\hookrightarrow M$.
\end{enumerate}
\hlt{Structural Commentary}: A split epimorphism is a quotient map with complemented kernel.\\
\hlt{Key Insight}: The right inverse chooses canonical representatives for the
cosets of $\ker\sigma$.
\end{proof}
\end{proposition}

\begin{principle} \label{princ:modules-not-vector-spaces} \hlt{Modules Are Not Just Vector Spaces with New Letters}\\
For a general ring $R$, the following vector-space implications fail:
\begin{enumerate}[label=(\roman*)]
\setlength{\itemsep}{0pt}
\item $N\leq M\nRightarrow\exists\,H\leq M$ with $M=N\oplus H$.
\item $M$ finitely generated and $N\leq M\nRightarrow N$ finitely generated.
\item $M\neq0\nRightarrow\exists\,\emptyset\neq S\subseteq M$ linearly independent.
\item $S$ minimal spanning or maximal independent $\nRightarrow S$ is a basis.
\item $M$ free and $N\leq M\nRightarrow N$ free.
\item $F$ free and $S\leq F\nRightarrow F/S$ free.
\end{enumerate}
Over a field, these obstructions vanish because $R\setminus\{0\}\subseteq R^\times$.
\end{principle}
